\subsection{一致收敛的一致结构}

设 X 是一个集，Y 是一致空间。对 Y 的每个围域 V，用 \(W(V)\) 表示 X 到 Y 中的映射的对 \((u, v)\) 中使 \((u(x), v(x)) \in V\)（对所有 \(x \in X\)）者组成的集。当 V 取遍 Y 的围域集时，诸集 \(W(V)\) 构成 F(X; Y) 上某个一致结构的围域基本系。因为它们显然满足公理 \((\mathrm{U}_{\mathrm{I}}^{\prime})\)（第 II 章，§ 1，第 1 号）；若 V 与 \(V'\) 是 Y 的两个使 \(V \subset V'\) 的围域，我们有 \(W(V) \subset W(V')\)，因而诸集 \(W(V)\) 满足 \((\mathrm{B}_{\mathrm{I}})\)（第 I 章，§ 6，第 3 号）；我们有 \(\widehat{W(V)} = W(\overline{V}^{-1})\)，故 \((\mathrm{U}_{\mathrm{II}}^{\prime})\) 满足；最后，关系式`` \((u(x), v(x)) \in V\)（对所有 \(x \in X\)）''与`` \((v(x), w(x)) \in V\)（对所有 \(x \in X\)）''蕴含关系式`` \((u(x), w(x)) \in V^{2}\)（对所有 \(x \in X\)）''；换言之，我们有 \(\widehat{W(V)} \subset W(\overline{V}^{2})\)，这就证明了 \((\mathrm{U}_{\mathrm{III}}^{\prime})\)。

定义 1 。在集 \(\mathcal{F}(\mathbf{X};\mathbf{Y})\) 上以子集 \(W(\mathbf{V})\)（其中 \(\mathbf{V}\) 取遍 \(\mathbf{Y}\) 的围域集）组成的集为围域基本系的一致结构，称为一致收敛的一致结构。它诱导的拓扑称为一致收敛的拓扑。若滤子 \(\Phi\)（\(\mathcal{F}(\mathbf{X};\mathbf{Y})\) 上的）关于这拓扑收敛于元素 \(u_0\)，则称 \(\Phi\) 一致收敛于 \(u_0\)。

注意 \(\mathcal{F}(\mathrm{X};\mathrm{Y})\) 上的一致收敛拓扑依赖于 Y 的一致结构，而不只是 Y 的拓扑（习题 4）。

把 \(\mathcal{F}(\mathbf{X};\mathbf{Y})\) 赋以一致收敛的一致结构所得的一致空间记作 \(\mathcal{F}_u(\mathbf{X};\mathbf{Y})\)。

